\documentclass[11pt,a4paper]{article}
\usepackage[utf8]{inputenc}
\usepackage[T1]{fontenc}
\usepackage{amsmath,amssymb,amsthm}
\usepackage{listings}
\usepackage{geometry}
\usepackage{hyperref}
\geometry{margin=2.5cm}

\newtheorem{theorem}{Theorem}
\newtheorem{lemma}[theorem]{Lemma}

\lstset{
  language=C++,
  basicstyle=\ttfamily\small,
  keywordstyle=\bfseries,
  breaklines=true,
  numbers=left,
  numberstyle=\tiny,
  frame=single,
  tabsize=2
}

\title{IOI 2022 --- Digital Circuit}
\author{}
\date{}

\begin{document}
\maketitle

\section{Problem Statement}

A digital circuit has $N$ \emph{threshold gates} (numbered $0,\ldots,N{-}1$)
and $M$ \emph{source gates} (numbered $N,\ldots,N{+}M{-}1$).  Gate~0 is the
output.  Each non-root gate has a unique parent (a threshold gate).  Every
threshold gate $i$ with $c_i$ inputs has a parameter $p_i \in [1, c_i]$: the
gate outputs 1 iff at least $p_i$ of its inputs are 1.

Source gates have binary states.  After each \emph{update} (toggling source
gates in a contiguous range $[L, R]$), count the number of parameter
assignments that make gate~0 output~1, modulo $10^9 + 22$.

\section{Solution}

\subsection{Per-gate counting}

For each threshold gate $i$ with $c_i$ inputs, if exactly $k$ of its inputs
output~1, then exactly $\min(k, c_i)$ choices of $p_i$ make gate~$i$
output~1, and $c_i - \min(k, c_i)$ make it output~0.

\subsection{Aggregation with generating polynomials}

For each gate $i$, define:
\begin{align*}
  \mathrm{ways\_on}[i]  &= \text{number of parameter assignments in the
                            subtree of } i \text{ making } i \text{ output 1,}\\
  \mathrm{ways\_off}[i] &= \text{the complementary count.}
\end{align*}

For a source gate $s$: $\mathrm{ways\_on}[s] = \text{state}(s)$ and
$\mathrm{ways\_off}[s] = 1 - \text{state}(s)$.

For a threshold gate $i$ with children $g_1, \ldots, g_{c_i}$, we accumulate
two quantities over the children:
\begin{align*}
  A &= \prod_j \bigl(\mathrm{ways\_on}[g_j] + \mathrm{ways\_off}[g_j]\bigr),\\
  B &= \sum_{\text{subsets }S}
       |S| \cdot \prod_{j \in S} \mathrm{ways\_on}[g_j]
             \cdot \prod_{j \notin S} \mathrm{ways\_off}[g_j].
\end{align*}
Then $\mathrm{ways\_on}[i] = B$ and
$\mathrm{ways\_off}[i] = c_i \cdot A - B$.

The pair $(A, B)$ can be maintained incrementally.  Writing
$\mathrm{tot}_j = \mathrm{ways\_on}[g_j] + \mathrm{ways\_off}[g_j]$,
the transition when appending child $g_{j+1}$ is
\begin{align*}
  A' &= A \cdot \mathrm{tot}_{j+1},\\
  B' &= B \cdot \mathrm{tot}_{j+1} + A \cdot \mathrm{ways\_on}[g_{j+1}].
\end{align*}

\subsection{Handling toggle updates}

When source gates in $[L,R]$ are toggled, the values
$\mathrm{ways\_on}$ and $\mathrm{ways\_off}$ change for all ancestors.
A full bottom-up recomputation takes $O(N + M)$ per update.

\begin{lemma}
The total work per update is $O(N + M)$ since each gate is recomputed
at most once in the bottom-up pass.
\end{lemma}

For full-score performance, one can factor each gate's contribution
as a product of per-source-gate affine functions and maintain these
products under range toggles using a segment tree with lazy
propagation.  We present the simpler $O((N+M)Q)$ recomputation below.

\section{Implementation}

\begin{lstlisting}
#include <bits/stdc++.h>
using namespace std;
typedef long long ll;
const int MOD = 1000002022;

int N, M;
vector<int> par;
vector<vector<int>> ch;
vector<int> src_state;

ll ways_on_0;

void recompute() {
    vector<ll> won(N + M), woff(N + M);
    for (int i = N; i < N + M; i++) {
        won[i]  = src_state[i - N];
        woff[i] = 1 - src_state[i - N];
    }
    for (int i = N - 1; i >= 0; i--) {
        ll A = 1, B = 0;
        int ci = (int)ch[i].size();
        for (int g : ch[i]) {
            ll tot = (won[g] + woff[g]) % MOD;
            ll nA = A * tot % MOD;
            ll nB = (B % MOD * tot % MOD + A % MOD * (won[g] % MOD)) % MOD;
            A = nA;
            B = nB;
        }
        won[i]  = B % MOD;
        woff[i] = ((ll)ci % MOD * (A % MOD) % MOD - B % MOD + MOD) % MOD;
    }
    ways_on_0 = won[0];
}

void init(int _N, int _M, vector<int> P, vector<int> A) {
    N = _N; M = _M;
    par = P;
    src_state = A;
    ch.resize(N + M);
    for (int i = 1; i < N + M; i++)
        ch[P[i]].push_back(i);
}

int count_ways(int L, int R) {
    for (int i = L; i <= R; i++)
        src_state[i - N] ^= 1;
    recompute();
    return (int)(ways_on_0 % MOD);
}
\end{lstlisting}

\section{Complexity}

\begin{itemize}
  \item \textbf{Initialization:} $O(N + M)$.
  \item \textbf{Per update (above code):} $O(N + M)$ for full recomputation.
  \item \textbf{Per update (optimized):} $O(N + M + \log M)$ using a segment
        tree over source gates that maintains the product of affine
        contributions and supports range-toggle with lazy propagation.
  \item \textbf{Space:} $O(N + M)$.
\end{itemize}

\end{document}
